% =============================================================================
\section{结论}
\label{sec:conclusion}
本文对 $\alpha$-稳定 \Levy\ 噪声驱动的空间耦合意见动力学模型中的随机共振
给出了分析与数值相结合的研究，并把它用于理解公众对误信息的易感性。在解析
方面，我们建立了适定性与矩估计，证明了小噪声极限下到确定性动力学的均方收敛，
导出了空间分数阶 Fokker--Planck 方程与非微扰矩层级，给出了有序--无序相变的
平均场判据（$V_c=a/3b$），由 Kramers 逃逸率得到了最优噪声强度的无参数预言
$D^{*}=\Delta U/2$，并借助禁区定理证明了输入--输出互信息中非周期随机共振的
存在性。在数值方面，我们在 $50\times50$ 格点上确认了上述每一条预言，记录了
重尾驱动下共振的显著右移与展宽，通过首次穿越时间测量把该右移追溯到跳跃介导
的（非激活的）势垒跨越，并识别了误信息脆弱性窗口以及在该窗口内辟谣的非线性
回报。

统一的信息是双面的。\emph{对计算社会科学而言}，结果表明公众对错误信念的
易感性由使微弱真实信号可被检测的同一共振机制所支配；重尾的信息环境比高斯
环境更持久地（尽管不那么尖锐地）易感；而一个高度同步的公共领域的表面健康
不应被照单全收——同步性本身不携带任何认识论上的保证。\emph{对概率与统计
物理而言}，结果表明经典 Kramers/两态随机共振理论在高维、空间耦合、带跳跃
噪声的系统中存活下来了，前提是逃逸率要被替换为其重尾对应物；共振是稳健的，
但决定其位置的平衡条件并不稳健。

自然的扩展包括：把规则格点替换为经验的或无标度拓扑，纳入顽固者与异质主体，
允许真相偏置为有限的注意力而竞争，以及用真实的级联数据校准
$(\alpha,D,C)$。这些每一步都把模型从机制性的思想实验推向一个经过校准的
工具——一个可以用来回答"这个平台现在有多脆弱、干预应该投在哪里"的工具。
而本文建立的理论结果——适定性、小噪声极限与非周期共振的存在性——无论拓扑
如何都仍然成立。
